\section{Introduction}
\label{sec:introduction}
In this paper, we study the structured nonsmooth nonconvex composite problem
\begin{align}
    \min_{\bx\in\RR^n}\quad & F(\bx) := f(\bx) + h(\calA(\bx)),
    \tag{\textbf{P}}
    \label{prob:main}
\end{align}
where $f\colon\RR^n\to\RR$ is a possibly nonconvex function with a Lipschitz
continuous gradient, $\calA\colon\RR^n\to\RR^m$ is a smooth and possibly nonlinear
mapping, and $h\colon\RR^m\to\RR\cup\{+\infty\}$ is a proper closed convex function
whose proximal mapping is easy to evaluate; the complete assumptions are stated in
Section~\ref{sec:assumptions}. The two components play complementary roles: $f$ is
usually a smooth loss or data-fidelity term, whereas $h$ is a nonsmooth convex function
that imposes robustness, sparsity, or other structure on the transformed variable
$\calA(\bx)$
\cite{charisopoulos2021low,markovsky2008structured,parekh2015convex}. Several models
of practical interest take this form. 

\subsection{Motivating Examples}

\textbf{Robust phase retrieval}~\cite{duchi2019solving}.
Given measurement vectors $\mathbf a_i\in\RR^n$ and observations
$b_i\approx\langle\mathbf a_i,\bx^\star\rangle^2$, a fraction of which may be grossly
corrupted, one solves
\begin{equation*}
  \min_{\bx\in\RR^n}\;\frac1m\sum_{i=1}^m
  \bigl|\langle\mathbf a_i,\bx\rangle^2-b_i\bigr|.
\end{equation*}
This is~Problem \eqref{prob:main} with $f\equiv0$, $h=\frac1m\|\cdot\|_1$, and
$\calA(\bx)_i=\langle\mathbf a_i,\bx\rangle^2-b_i$; the $\ell_1$ loss makes the
model robust to outliers.

\textbf{Robust low-rank matrix recovery}~\cite{charisopoulos2021low}.
Given a linear measurement map $\mathcal M\colon\RR^{d_1\times d_2}\to\RR^m$ and
observations $\mathbf b$ containing outliers, a rank-$r$ matrix is recovered in
factorized form via
\begin{equation*}
  \min_{\mathbf U\in\RR^{d_1\times r},\,\mathbf V\in\RR^{d_2\times r}}
  \frac1m\bigl\|\mathcal M(\mathbf U\mathbf V^\top)-\mathbf b\bigr\|_1.
\end{equation*}
Here $\bx=(\mathbf U,\mathbf V)$, $h=\frac1m\|\cdot\|_1$, and
$\calA(\bx)=\mathcal M(\mathbf U\mathbf V^\top)-\mathbf b$ is bilinear in
$(\mathbf U,\mathbf V)$, so the model again has the form~Problem \eqref{prob:main}.

\textbf{Minimizing a pointwise maximum}~\cite{zhang2020single}.
For smooth losses $f_1,\dots,f_m$, such as the losses on $m$ data subgroups, one solves
\begin{equation*}
  \min_{\bx\in\RR^n}\;\max_{1\le i\le m} f_i(\bx).
\end{equation*}
This is~Problem \eqref{prob:main} with $h(u)=\max_{1\le i\le m}u_i$ and
$\calA(\bx)=(f_1(\bx),\dots,f_m(\bx))$; a smooth regularizer, if present, enters
through $f$, and the proximal mapping of $h$ reduces to a projection onto the unit
simplex.

\subsection{Related Work}

Introducing the auxiliary variable $\by=\calA(\bx)$ rewrites~Problem \eqref{prob:main} as
the equality-constrained problem
\begin{equation}\label{prob:split}
  \min_{\bx,\by}\; f(\bx)+h(\by)\quad\text{s.t.}\quad \calA(\bx)=\by,
\end{equation}
which brings it within the scope of penalty methods and augmented Lagrangian methods
(ALMs). For general nonlinearly constrained nonconvex problems, both families control
feasibility through the penalty parameter. Quadratic penalty methods increase the
penalty without bound~\cite{kong2019complexity,lin2022complexity}. ALMs were
originally introduced to avoid such unbounded penalties
\cite{hestenes1969multiplier,powell1969method}, and in convex programming any fixed
penalty suffices~\cite{rockafellar1973multiplier}; in the nonconvex setting, however,
inexact ALMs still increase the penalty across outer iterations to obtain complexity
guarantees~\cite{sahin2019inexact,li2021rate}, and proximal ALM, which keeps the
penalty fixed within a run, chooses its value according to the target accuracy
\cite{xie2021proximal}. Single-loop methods have attracted growing attention because
they avoid solving subproblems to high accuracy, yet they rely on the same mechanism:
single-loop quadratic penalty schemes increase the penalty~\cite{alacaoglu2024complexity},
and linearized ALMs use small dual steps oraccuracy $eps$-dependent parameters
\cite{shi2025momentum}. A large penalty inflates the smoothness constant of the
primal update, which forces small step sizes and makes the primal subproblems
increasingly ill-conditioned.

The split formulation~\eqref{prob:split} also carries favorable structure. Its
constraint Jacobian $[J_{\calA}(\bx),\,-I]$ has full row rank at every $\bx$, so
constraint qualifications hold automatically, and its objective separates into a
smooth part in $\bx$ and a prox-friendly part in $\by$. Nonconvex ADMM and related
Lagrangian-based splitting methods exploit exactly this structure
\cite{bot2020proximal,cohen2022dynamic,hallak2023adaptive,elbourkhissi2025perturbed}.
The structure, however, does not remove the need for a large penalty. For nonlinear
$\calA$, adaptive Lagrangian schemes require penalty levels proportional to
$1/\epsilon$~\cite{hallak2023adaptive}, and linearized perturbed ALMs use
accuracy-dependent penalty and perturbation parameters~\cite{elbourkhissi2025perturbed}.
The nonconvex--linear minimax route attains $\mathcal O(\epsilon^{-3})$
\cite{pan2021efficient,he2024approximation,xu2026rada}, but its dual regularization
acts as a penalty parameter that grows as the target accuracy decreases. Fixed-penalty
guarantees are available for smooth nonlinear equality constraints under a uniform LICQ
\cite{elbourkhissi2025linearized,liu2026fixed}, but they do not cover the nonsmooth
term $h(\by)$ in~\eqref{prob:split}. It therefore remains open whether a single-loop
ALM can keep every penalty-type parameter independent of both the accuracy
and the target accuracy while retaining an explicit first-order update for
$h\circ\calA$.

From another perspective, Problem~\eqref{prob:main} also admits a minimax reformulation. In particular, eliminating $\by$ from the Lagrangian of~\eqref{prob:split} gives
\begin{equation}\label{eq:ncl}
  \min_{\bx\in\RR^n}\max_{\blmbd\in\RR^m}\;
  \calL(\bx,\blmbd) =  \Phi(\bx,\blmbd)-h^*(\blmbd), \qquad \Phi(\bx,\blmbd):=f(\bx)+\langle\blmbd,\calA(\bx)\rangle,
\end{equation}
where $\operatorname{dom}h^*\subseteq\mathbb B(\mathbf 0,L_h)$ by
Assumption~\ref{ass:h}. This is a nonconvex--linear instance of the smooth
nonconvex--concave (NC-C) minimax problem. Since $\blmbd$ is exactly the multiplier of $\calA(\bx)=\by$, KKT points
of~\eqref{prob:split} and game-stationary (GS) points of~\eqref{eq:ncl} agree up to
constants, and ALMs can be read as single-loop minimax methods. 
For large-scale problems, we focus on single-loop first-order methods, which
update each variable by one explicit step per iteration without an inner
solver. The basic method of this kind for~\eqref{eq:ncl} is gradient
descent--ascent (GDA). It converges slowly because regularity fails on both
sides: $\Phi$ is only linear in $\blmbd$, so the best response
$\blmbd^\star(\bx)$ is in general set-valued, and $\Phi$ is nonconvex in $\bx$,
so $\bx^\star(\blmbd)$ need not be unique or stable. Even with two timescales,
GDA needs $\mathcal O(\epsilon^{-6})$ iterations to reach an $\epsilon$-GS
point~\cite{lin2020gradient}.

Two techniques improve on this rate, each restoring regularity on one side.
The first is primal smoothing, which adds $\frac{1}{2\gamma}\|\bx-\bz\|^2$ with
a fixed $\gamma$ and moves the center slowly by
$\bz^{t+1}=\bz^t+\beta(\bx^{t+1}-\bz^t)$. The target accuracy enters only
through $\beta$: with $\beta=\Theta(\epsilon^2)$, the resulting Smoothed GDA
finds an $\epsilon$-GS point in $\mathcal O(\epsilon^{-4})$ iterations. For the
pointwise maximum of finitely many smooth functions, i.e.,~\eqref{eq:ncl} with
$h(\bu)=\max_i u_i$, strict complementarity permits a constant $\beta$ and
improves the bound to $\mathcal O(\epsilon^{-2})$~\cite{zhang2020single}.

The second technique is dual regularization, which subtracts
$\frac{1}{2s}\|\blmbd-\bmu\|^2$ from the dual objective and makes it strongly
concave. In our reduced augmented-Lagrangian formulation, this regularization
has the following combined penalty interpretation:
\begin{equation}\label{eq:intro-effective-penalty}
  \max_{\blmbd}\mathcal P_s(\bx,\blmbd;\bz,\bmu)
  =\calD_{\rho+s}(\bx,\bmu)
    +\frac{1}{2\gamma}\|\bx-\bz\|^2,
\end{equation}
where $\mathcal P_s$ is defined in~\eqref{eq:potential} and
$\calD_{\rho+s}$ is the reduced augmented Lagrangian~\eqref{eq:reduced-al}
with the combined penalty parameter $\rho+s$; see Lemma~\ref{lem:effective-penalty}.
A zero center $\bmu=\mathbf 0$ recovers a quadratic penalty with
coefficient $\rho+s$, which can leave a nonzero feasibility residual
when $s$ is fixed. Updating $\bmu$ instead provides a multiplier-update
mechanism that may recover feasibility while keeping $\rho+s$ fixed;
its convergence and complexity depend on the coupled updates.

With the dual center fixed at $\bmu=\mathbf 0$, Perturbed GDA
attains an $\mathcal O(\epsilon^{-4})$ complexity for finding an
$\epsilon$-GS point~\cite{lu2020hybrid,xu2023unified}, and Li et al. establish
a matching algorithm-specific lower bound~\cite{li2026smoothing}.
Methods exploiting the NC-L structure attain
$\mathcal O(\epsilon^{-3})$~\cite{pan2021efficient,he2024approximation,xu2026rada}.
These guarantees rely on accuracy-dependent regularization parameters,
corresponding to a penalty parameter that increases as $\epsilon$ decreases. To improve the iteration complexity, Li et al. combine primal smoothing
with zero-centered dual perturbation to develop Perturbed Smoothed GDA,
which yields a genuine improvement in complexity order, reducing
the GS complexity of Perturbed GDA from $\mathcal O(\epsilon^{-4})$
to $\mathcal O(\epsilon^{-3})$ with a matching lower
bound~\cite{li2026smoothing}. However, the dual regularization coefficient still satisfies
$1/s=\Theta(\epsilon)$.

With a moving dual center $\bmu$ and fixed primal and dual regularization
coefficients, DS-GDA achieves $\mathcal O(\epsilon^{-4})$ complexity
for nonconvex--concave problems. Faster rates are available under a
uniform one-sided K{\L} condition~\cite{zheng2023universal}.
Without primal smoothing, PPALA reports
$\widetilde{\mathcal O}(\epsilon^{-2})$ complexity with fixed parameters
for functionally constrained optimization~\cite{kim2024fast}.
However, its feasibility argument relies on asymptotic reasoning,
and a corresponding finite-iteration feasibility bound remains
to be verified.

In all guarantees above that reach $\mathcal O(\epsilon^{-3})$, some
parameter grows with $1/\epsilon$: either the penalty, as in penalty methods
and GDPA, or the dual regularization, as in Perturbed Smoothed GDA. Practical
ALM implementations, in contrast, usually keep the penalty fixed. Growing the
penalty is costly for a single-loop method, because the Lipschitz constant of
$\nabla_{\bx}\calD_\rho$ increases with $\rho$ and forces the primal stepsize
to shrink. This leads to the following question.
\begin{quote}
  \textbf{(Q)}
  \emph{To reach $\mathcal O(\epsilon^{-3})$, must a single-loop linearized ALM
  let some parameter grow with $1/\epsilon$? If so, must this growth be
  carried by the penalty parameter?}
\end{quote}
In this paper, we provide a precise characterization addressing (Q). We study
SPLALM, a family of single-loop linearized ALMs that perform one gradient
descent--ascent step per iteration on $\calD_\rho$ with a fixed penalty
$\rho$, a dual regularization with parameter $s$, and moving primal and dual
centers (Definition~\ref{def:splalm}). Every member of the family has the
combined penalty $\rho+s$ (Lemma~\ref{lem:effective-penalty}). We answer the
first part of (Q) affirmatively and the second part negatively. If
$\rho+s\le c_1/\epsilon$, no choice of constant center stepsizes improves on
$\Omega(\epsilon^{-4})$ (Theorem~\ref{thm:lower-bound}). If $\rho$ is fixed
and $s$ is of order $1/\epsilon$, suitable center stepsizes give
$\mathcal O(\epsilon^{-3})$ (Theorem~\ref{thm:upper-beta0}). Hence some
parameter must grow, and the dual regularization alone can carry this growth.
The lower bound is new. The upper bounds adapt the analyses of Smoothed GDA
and Perturbed Smoothed GDA, both members of the family, thus demonstrating the
tightness of the lower bound and of the critical order of the combined
penalty. Table~\ref{tab:known-results} compares these results with known
guarantees. Our contribution can be summarized as follows: 
\begin{table*}[tp]
\centering
\begin{threeparttable}
\footnotesize
\setlength{\tabcolsep}{6pt}
\renewcommand{\arraystretch}{1.15}
\begin{tabular}{@{}lccc@{}}
\toprule
Method & Parameter scale & Upper bound & Lower bound \\
\midrule
\makecell[l]{Smoothed GDA\\\cite{zhang2020smoothed,li2026smoothing}}
 & Fixed primal smoothing
 & $\mathcal O(\epsilon^{-4})$ & $\Omega(\epsilon^{-4})$ \\[3pt]
\makecell[l]{Perturbed Smoothed GDA\\\cite{li2026smoothing}}
 & $s=\Theta(\epsilon^{-1})$
 & $\mathcal O(\epsilon^{-3})$ & $\Omega(\epsilon^{-3})$ \\[3pt]
% \midrule 
\makecell[l]{GDPA\\\cite{luSingleLoopGradientDescent2022}}
 & $\rho_k=\Theta(\epsilon^{-1})$
 & $\mathcal O(\epsilon^{-3})$ & --- \\[3pt]
\midrule
\rowcolor{gray!15}
\makecell[l]{\textbf{SPLALM family}\\Theorems~\ref{thm:upper-beta1} and~\ref{thm:lower-bound}}
 & \makecell{$\rho=\Theta(1)$,\\$\rho+s\le c_1\epsilon^{-1}$}
 & $\mathcal O(\epsilon^{-4})$ & $\Omega(\epsilon^{-4})$ \\[3pt]
\makecell[l]{\textbf{SPLALM family}\\Theorem~\ref{thm:upper-beta0}}
 & \makecell{$\rho=\Theta(1)$,\\$s=c_s\epsilon^{-1}$}
 & $\mathcal O(\epsilon^{-3})$ & --- \\[3pt]
\bottomrule
\end{tabular}
\caption{Complexity guarantees for representative single-loop methods.
For SPLALM, the bounds concern the family complexity in
Definition~\ref{def:complexity-measure} over
$\mathcal I(\mathbf C_0,\Delta)$, under the conditions of the cited theorems.
In the last row, $c_s$ is chosen as in Theorem~\ref{thm:upper-beta0}.
The rates use each method's stated stationarity criterion and assumptions;
lower bounds are algorithm-specific. ``---'' denotes an unreported bound.}
\label{tab:known-results}
\end{threeparttable}
\end{table*}
\begin{enumerate}
    \item \textbf{A lower bound independent of the center stepsizes.}
    We construct hard instances on which
    SPLALM, for every choice of constant center stepsizes
    $(\alpha,\beta)\in[0,1]^2$, requires $\Omega(\epsilon^{-4})$ iterations
    whenever $\epsilon(\rho+s)\le c_1$ (Theorem~\ref{thm:lower-bound}).
    Moving the centers therefore cannot compensate for a combined penalty
    $\rho+s=o(\epsilon^{-1})$.
    % TODO: GS version. A strictly convex perturbation of the outer function
    % would convert this bound into a lower bound of the same order for game
    % stationarity of the output pairs (x^{k+1},q^k). Restore this sentence
    % only after the perturbation lemma is proved.

    \item \textbf{Tightness for a fixed penalty.}
    For $\beta=1$, SPLALM reduces to Smoothed GDA applied to $\calD_\rho$ and
    attains $\mathcal O(\epsilon^{-4})$ uniformly over
    $\mathcal I(\mathbf C,\Delta)$ (Theorem~\ref{thm:upper-beta1}). Hence the
    lower bound is tight when $\rho$ is fixed and $\epsilon(\rho+s)\le c_1$,
    for sufficiently small $\gamma$, $\eta$, and $\tau$.

    \item \textbf{The critical scale of the combined penalty.}
    For $\beta=0$, $\rho$ fixed, and $s=c_s/\epsilon$, SPLALM reduces to
    Perturbed Smoothed GDA applied to $\calD_\rho$ and attains
    $\mathcal O(\epsilon^{-3})$ (Theorem~\ref{thm:upper-beta0}). Together with
    the lower bound, $\epsilon^{-1}$ is the critical order of $\rho+s$, and
    this growth can be carried by $s$ alone
    (Corollary~\ref{cor:effective-penalty}).
\end{enumerate}

% Open questions moved from the former (Q2) and (Q3); to be placed in the
% concluding section: (i) the complexity gain of primal smoothing under a
% fixed s and a moving dual center; (ii) structural conditions on
% (f,h,A) that yield a center--primal error bound and faster rates under a
% fixed combined penalty.

\paragraph{Organization}
Section~\ref{sec:preliminaries} introduces the problem class and the KKT and
game-stationarity criteria. Section~\ref{sec:algorithm-framework} introduces
the SPLALM family and the complexity measure, and shows that every member
of the family has the combined penalty
$\rho+s$. Section~\ref{sec:lower-bound} states the main results, and
Section~\ref{sec:lb-proof} proves the lower bound.
Section~\ref{sec:conclusion} discusses open questions.
